% Dokumenteinstellungen und Anpassungen
\documentclass[chapterprefix=false, 12pt, a4paper, oneside, parskip=half, listof=totoc, bibliography=totoc, numbers=noendperiod]{article}

%\renewcommand{\familydefault}{\sfdefault}
%\usepackage{helvet}


%Anpassung der Seitenränder (Standard bottom ca. 52mm abzüglich von ca. 4mm für die nach oben rechts gewanderte Seitenzahl)
\usepackage[bottom=20mm,left=25mm,right=25mm, top=25mm]{geometry}

\usepackage{csquotes}
\usepackage[round,sort&compress, authoryear]{natbib}
\bibliographystyle{plainnat}

%Tweaks für scrbook
\usepackage{scrhack}

\usepackage{tabularx} % in the preamble

%Blindtext
\usepackage{blindtext}

%Erlaubt unter anderem Umbrüche captions
\usepackage{caption}

%Stichwortverzeichnis
\usepackage{imakeidx}

%Kompakte Listen
\usepackage{paralist}

%Zitate besser formatieren und darstellen
\usepackage{epigraph}

%Glossar, Stichwortverzeichnis (Akronyme werden als eigene Liste aufgeführt)
\usepackage[toc, acronym]{glossaries}

%Anpassung von Kopf- und Fußzeile
%beeinflusst die erste Seite des Kapitels
\usepackage[automark,headsepline]{scrlayer-scrpage}
%\input{resources/styles/header_footer}
\setlength{\parindent}{0em}

%\renewcommand\chapter{\thispagestyle{plain}}

%Auskommentieren für die Verkleinerung des vertikalen Abstandes eines neuen Kapitels
%\renewcommand*{\chapterheadstartvskip}{\vspace*{.25\baselineskip}}
%\renewcommand*{\chapterheadendvskip}{\vspace*{.25\baselineskip}}

%Zeilenabstand 1,5
\usepackage[onehalfspacing]{setspace}

%Verbesserte Darstellung der Buchstaben zueinander
\usepackage[stretch=10]{microtype}

%Deutsche Bezeichnungen für angezeigte Namen (z.B. Inhaltsverzeichnis etc.)
\usepackage[ngerman]{babel}

%Unterstützung von Umlauten und anderen Sonderzeichen (UTF-8)
\usepackage{lmodern}
\usepackage[utf8]{luainputenc}
\usepackage[T1]{fontenc}
\usepackage[ngerman]{babel}
\setlength{\emergencystretch}{1em}

%Einfachere Zitate
\usepackage{epigraph}

%Unterstützung der H Positionierung (keine automatische Verschiebung eingefügter Elemente)
\usepackage{float}

%Erlaubt Umbrüche innerhalb von Tabellen
\usepackage{tabularx}

%Erlaubt Seitenumbrüche innerhalb von Tabellen
\usepackage{longtable}

%Erlaubt die Darstellung von Sourcecode mit Highlighting
\usepackage{listings}

%Definition eigener Farben bei Nutzung eines selbst vergebene Namens
\usepackage[table,xcdraw]{xcolor}

%Vektorgrafiken
\usepackage{tikz}

%Grafiken (wie jpg, png, etc.)
\usepackage{graphicx}

%Grafiken von Text umlaufen lassen
\usepackage{wrapfig}

%Ermöglicht Verknüpfungen innerhalb des Dokumentes (e.g. for PDF), Links werden durch "hidelink" nicht explizit hervorgehoben
\usepackage[hidelinks,german]{hyperref}

\setlength{\headheight}{18pt}

%Zusätzliche Farben
\definecolor{darkgreen}{RGB}{0,100,0}
\setlength{\emergencystretch}{15pt}

%Umbenennungen
\renewcommand{\lstlistlistingname}{Quelltextverzeichnis}

\renewcommand{\lstlistingname}{Codeauszug}
\lstset{
	language=Java,
	numbers=left,
	columns=fullflexible,
	aboveskip=5pt,
	belowskip=10pt,
	basicstyle=\small\ttfamily,
	backgroundcolor=\color{black!5},
	commentstyle=\color{darkgreen},
	keywordstyle=\color{blue},
	stringstyle=\color{gray},
	showspaces=false,
	showstringspaces=false,
	showtabs=false,
	xleftmargin=16pt,
	xrightmargin=0pt,
	framesep=5pt,
	framerule=3pt,
	frame=leftline,
	rulecolor=\color{green},
	tabsize=2,
	breaklines=true,
	breakatwhitespace=true,
	prebreak={\mbox{$\hookleftarrow$}}
}

%Anpassungen für das Abkürzungsverzeichnis
\newglossarystyle{dottedlocations}{%
	\renewcommand*{\glossaryentryfield}[5]{%
		\item[\glsentryitem{##1}\glstarget{##1}{##2}] \emph{##3}%
		\unskip\leaders\hbox to 2.9mm{\hss.}\hfill##5}%
	\renewcommand*{\glsgroupskip}{}%
}


\usepackage{graphicx}
\usepackage{subcaption}
\usepackage{float}
\usepackage{color}
\usepackage{listings}
\usepackage{eso-pic}
\usepackage{caption}
\newcommand\BackgroundPic{%
\put(0,0){%
\parbox[b][\paperheight]{\paperwidth}{%
\vfill
\centering
\includegraphics[width=\paperwidth,height=\paperheight,%
keepaspectratio]{resources/images/deckblatt.jpg}%
\vfill
}}}


\makeatletter
\newcommand*{\germanTitle}[1]{\gdef\@germanTitle{#1}}
\newcommand*{\englishTitle}[1]{\gdef\@englishTitle{#1}}
\newcommand*{\gradeType}[1]{\gdef\@gradeType{#1}}
\newcommand*{\firstExaminer}[1]{\gdef\@firstExaminer{#1}}
\newcommand*{\secondExaminer}[1]{\gdef\@secondExaminer{#1}}
\newcommand*{\matrikelnr}[1]{\gdef\@matrikelnr{#1}}
\newcommand*{\authorBirthplace}[1]{\gdef\@authorBirthplace{#1}}
\newcommand*{\submitDate}[1]{\gdef\@submitDate{#1}}
\newcommand*{\discipline}[1]{\gdef\@discipline{#1}}
\newcommand*{\courseOfStudies}[1]{\gdef\@courseOfStudies{#1}}
\newcommand*{\authorLastname}[1]{\gdef\@authorLastname{#1}}
\newcommand*{\authorFirstname}[1]{\gdef\@authorFirstname{#1}}


\renewcommand*{\maketitle}{
	\begin{titlepage}
		\newgeometry{left=3.8cm,right=2.5cm,top=4.6cm,bottom=2.5cm}
			\begingroup
				\fontsize{44pt}{46pt}\selectfont
				{\bfseries Seminarfacharbeit}
			\endgroup

			\vskip 1.44cm

			\begingroup
			\fontsize{8pt}{6pt}\selectfont
			Titel der Arbeit // Title of Thesis
			\endgroup

			\vskip -0.02cm

			\begingroup
			\fontsize{16pt}{16pt}\selectfont
			{\bfseries \@germanTitle\par}
			{\bfseries \@englishTitle\par}
			\endgroup
			\vskip -0.1cm

			\noindent\rule{15cm}{0.4pt}
			
			\vskip 0.05cm

			\begingroup
			\fontsize{8pt}{6pt}\selectfont
			Akademischer Abschlussgrad: Grad, Fachrichtung (Abkürzung) // Degree
			\endgroup

			\vskip -0.15cm

			\begingroup
			\fontsize{12pt}{14pt}\selectfont
				{\@gradeType}
			\endgroup
			\vskip -0.15cm

			\noindent\rule{15cm}{0.4pt}
			
			\vskip 0.05cm

			\begingroup
			\fontsize{8pt}{6pt}\selectfont
			Autorenname, Geburtsort // Name, Place of Birth
			\endgroup

			\vskip -0.15cm

			\begingroup
			\fontsize{12pt}{14pt}\selectfont
				{\@authorFirstname} {\@authorLastname}, {\@authorBirthplace}
			\endgroup
			\vskip -0.15cm

			\noindent\rule{15cm}{0.4pt}
			
			\vskip 0.05cm

			\begingroup
			\fontsize{8pt}{6pt}\selectfont
			Studiengang // Course of Study
			\endgroup

			\vskip -0.15cm

			\begingroup
			\fontsize{12pt}{14pt}\selectfont
				{\@courseOfStudies}
			\endgroup
			\vskip -0.15cm

			\noindent\rule{15cm}{0.4pt}
			
			\vskip 0.05cm

			\begingroup
			\fontsize{8pt}{6pt}\selectfont
			Fachbereich // Department
			\endgroup

			\vskip -0.15cm

			\begingroup
			\fontsize{12pt}{14pt}\selectfont
				{\@discipline}
			\endgroup
			\vskip -0.15cm

			\noindent\rule{15cm}{0.4pt}
			
			\vskip 0.05cm

			\begingroup
			\fontsize{8pt}{6pt}\selectfont
			Erstprüferin/Erstprüfer // First Examiner
			\endgroup

			\vskip -0.15cm

			\begingroup
			\fontsize{12pt}{14pt}\selectfont
				{\@firstExaminer}
			\endgroup
			\vskip -0.15cm

			\noindent\rule{15cm}{0.4pt}
			
			\vskip 0.05cm

			\begingroup
			\fontsize{8pt}{6pt}\selectfont
			Zweitprüferin/Zweitprüfer // Second Examiner
			\endgroup

			\vskip -0.15cm

			\begingroup
			\fontsize{12pt}{14pt}\selectfont
				{\@secondExaminer}
			\endgroup
			\vskip -0.15cm

			\noindent\rule{15cm}{0.4pt}
			
			\vskip 0.05cm

			\begingroup
			\fontsize{8pt}{6pt}\selectfont
			Abgabedatum // Date of Submission
			\endgroup

			\vskip -0.15cm

			\begingroup
			\fontsize{12pt}{14pt}\selectfont
				{\@submitDate}
			\endgroup
			\vskip -0.15cm

			\noindent\rule{15cm}{0.4pt}
		\restoregeometry
	\end{titlepage}
}
\makeatother

% Variablen für das Deckblatt
\gradeType{Abitur 2021}
\germanTitle{Face Recognition mit Hilfe von Künstlicher Intelligenz}
\englishTitle{}
\authorFirstname{Kevin}
\authorLastname{Pagenkämper}
\authorBirthplace{Nordhorn}
\discipline{Technik und Ethik}
\courseOfStudies{Informatik}
\matrikelnr{112}
\submitDate{28.02.2020}
\firstExaminer{Andre Nixdorf}
\secondExaminer{Henrike Schnöing}


\makeatletter

\newcommand*{\place}[1]{\gdef\@place{#1}}

\newcommand*{\makeeidesstatt}{
	\begin{titlepage}
		\newgeometry{left=3.8cm,right=2.5cm,top=8cm,bottom=2.5cm}
			\begingroup
			\fontsize{18pt}{20pt}\selectfont
			{\bfseries Eidesstattliche Versicherung}
			\endgroup

			\vskip 0.8cm

			\begingroup
			\fontsize{12pt}{14pt}\selectfont
				{\@authorLastname} {\@authorFirstname}
			\endgroup

			\vskip -0.5cm

			\noindent\rule{15cm}{0.4pt}

			\vskip -0.4cm

			\begingroup
			\fontsize{8pt}{6pt}\selectfont
			Name, Vorname // Name, First Name
			\endgroup

			\vskip 0.6cm
			\begingroup
			\fontsize{10.5pt}{11.5pt}\selectfont
			Ich versichere hiermit an Eides statt, dass ich die vorliegende Seminarfacharbeit mit dem Titel
			\endgroup

			\begingroup
			\fontsize{16pt}{18pt}\selectfont
			{\bfseries \@germanTitle}
			\endgroup

			\begingroup
			\fontsize{10.5pt}{11.5pt}\selectfont
			selbstständig und ohne unzulässige fremde Hilfe erbracht habe. Ich habe keine anderen als die angegebenen Quellen und Hilfsmittel benutzt sowie wörtliche und sinngemäße Zitate kenntlich gemacht. Die Arbeit hat in gleicher oder ähnlicher Form noch keiner Prüfungsbehörde vorgelegen.
			\endgroup

			\vskip 0.8cm

			{
			\fontsize{12pt}{14pt}\selectfont
			{\@place}, den {\today}
			}

			\vskip -0.5cm

			\noindent\rule{15cm}{0.4pt}

			\vskip -0.3cm

			\begingroup
			\fontsize{8pt}{6pt}\selectfont
			Ort, Datum, Unterschrift // Place, Date, Signature
			\endgroup


		\restoregeometry
	\end{titlepage}
}
\makeatother

\place{Nordhorn}


\lstset{literate=%
  {Ö}{{\"O}}1
  {Ä}{{\"A}}1
  {Ü}{{\"U}}1
  {ß}{{\ss}}1
  {ü}{{\"u}}1
  {ä}{{\"a}}1
  {ö}{{\"o}}1
}

% Javascript als Sprache für die lstings
\lstdefinelanguage{JavaScript}{
  keywords={typeof, new, true, false, catch, function, return, null, catch, switch, var, if, in, while, do, else, case, break, that, globals, let, const},
  keywordstyle=\color{blue}\bfseries,
  ndkeywords={class, export, boolean, throw, implements, import, this},
  ndkeywordstyle=\color{darkgray}\bfseries,
  identifierstyle=\color{black},
  sensitive=false,
  comment=[l]{//},
  morecomment=[s]{/*}{*/},
  commentstyle=\color{gray}\ttfamily,
  stringstyle=\color{red}\ttfamily,
  morestring=[b]',
  morestring=[b]"
}
\lstset{
    language=JavaScript,
    numbers=none,
    frame=leftline,
    tabsize=2,
    rulesepcolor=\color{gray},
    rulecolor=\color{black},
    captionpos=b,
    breaklines=true,
    breakatwhitespace=false,
}

% Neuer Befehl um Subsubsubkapitel schreiben zu können
\newcommand{\subsubsubsection}[1]{\paragraph{#1}\mbox{}\\}
\setcounter{secnumdepth}{4}
\setcounter{tocdepth}{4}
\usepackage{fontspec}
\setmainfont{Arial}

\begin{document}

% Entfernen der Seitenzahlen und BackroundPic als Hintergrund nutzen
\pagenumbering{gobble}
\AddToShipoutPicture{\BackgroundPic}

% Titelblatt erzeugen
\maketitle
\newpage

% Eidesstattliche Erklärung erzeugen
\makeeidesstatt
\thispagestyle{empty}

% Abstract und Inhaltsverzeichnisse einfügen (Numerierung römisch ab Inhaltsverzeichnis)
\section*{Abstract}
\label{sec:abstract}

\ClearShipoutPicture
\thispagestyle{empty}
\newpage
\pagenumbering{Roman}
\tableofcontents
%\newpage
%\listoffigures
\newpage

% Kapitel einfügen
\pagenumbering{arabic}
% „ = Alt + 0132
% “ = Alt + 0147
\section{Einleitung}
\label{sec:einleitung}
Vor ungefähr 64 Jahren wurde das erste Programm, welches speziell zur Problemlösung und nachahmung eines Menschen entwickelt wurde, von Herbert Alexander Simon und Allen Newell geschrieben. Dieses Programm war Quellen zufolge das erste Programm welches der Definition von Künstlicher Intelligenz, ein Programm welches speziell dafür entwickelt wurde die Fähigkeiten eines Menschen Nachzuahmen, gerecht wurde [vgl. \ref{bib:LogicTheorist}] wobei zu dem Zeitpunkt der Entwicklung dieser das Wort "Künstliche Intelligenz" [folgend auch als KI bezeichnet] noch nicht existierte bzw. nicht in diesem Kontext verwendet wurde [vgl. \ref{LogicTheorist:Wikipedia}]. Seitdem hat sich das Thema KI sowohl in der Hardware aber auch in der Software stark weiterentwickelt. Wo Künstliche Intelligenzen früher nur mit einfachen Zahlen arbeiten konnten, können einige mittlerweile Personen, Gesichter, Objekte und andere Dinge identifizieren und differenzieren. Doch auch hierbei gibt es Probleme die seit der Entwicklung von Gesichtserkennungsprogrammen in den 1960er nicht vollständig gelöst werden konnten, weshalb es bisher auch keine KI gibt welche keine Mängel aufweist. Und gibt es eine KI die momentan eine erfolgschance von 100\% hat ist dies aus dem Grund so, dass sie ihrem Schwachpunkt noch nicht begegnet ist. In dieser Facharbeit wird das Thema der Gesichtserkennung, zu engl. "Face recognition", behandelt. Bei dieser Art der Anwendung von Künstlicher Intelligenz liegt die Schwierigkeit besonders bei Unschärfe, Fragmenten, Verzerrungen und Ähnlichkeiten von Objekten und besonders von Gesichtern. Dies ist z.B. bei Zwillingen der Fall an dem viele Systeme noch heute scheitern. Zudem wird die Theorie von Face Recognition mit Hilfe von Künstlicher Intelligenz erklärt, mit praktischen Beispielen erläutert und auf ethischer Basis analysiert.
\newpage
\section{Literatur Verzeichnis}
\label{sec:literatur}

\subsection{Logic Theorist}
\label{bib:LogicTheorist}
	Letzter Zugriff: 12.02.2020\\
	URL: https://history-computer.com/ModernComputer/Software/LogicTheorist.html

	\subsubsection{Memorandum}
	\label{LogicTheorist:Memorandum}
		Letzter Zugriff: 18.02.2020 12:38\\
		URL: https://history-computer.com/Library/Logic\%20Theorist\%20memorandum.pdf
		
	\subsubsection{Wikipedia Artikel}
	\label{LogicTheorist:Wikipedia}
	Letzter Zugriff: 18.02.2020\\
	URL: https://en.wikipedia.org/wiki/Logic\_Theorist

\subsection{Allgemeine Quellen zum Thema KI}
\label{bib:AllgeimeineQuellen}
		\subsubsection{Wikipedia Artikel zum Thema Künstliche Intelligenz}
		\label{wiki:KuenstlicheIntelligenz}
		Letzter Zugriff: 19.02.2020 12:25\\
		URL: https://de.wikipedia.org/wiki/Künstliche\_Intelligenz

		\subsubsection{Gesichtserkennung}
		\label{wiki:face_recognition}
		Letzter Zugriff: 02.03.2020 17:19\\
		URL: https://de.wikipedia.org/wiki/Gesichtserkennung
	
\subsection{Aufbau und Funktionsweise Neuronaler Netze}		
\label{subsec:Aufbau_Funktion_Neuronaler_Netze}
		\subsubsection{JAAI - Aufbau von Neuronalen Netzen}
		\label{subsubsec:Aufbau_von_Neuronalen_Netzen}
		Letzter Zugriff: 22.02.22 15:04\\
		URL: https://jaai.de/kuenstliche-neuronale-netze-aufbau-funktion-291/

		\subsubsection{Wie lernen Neuronale Netze?}
		\label{subsubsec:Wie_lernen_neuronale_netze}
		Letzter Zugriff: 23.02.2020 17:49\\
		URL: https://jaai.de/machine-deep-learning-529/

\subsection{Bücher}
\label{books}

	\subsubsection{Buch: Digitale Gesichtserkennung}
	\label{book:DigitalGesichtserkennung}
		Untertitel: Theoretischer Überblick und praktische C++-Implementierung\\
		Autor: Andreas G. Ranftl\\
		Jahr: 2012\\
		Ort: Hamburg\\
		ISBN: 978-3-86341-432-0\\

	\subsubsection{Buch: Künstliche Intelligenz - Wann übernehmen die Maschinen?}
	\label{book:KI_WannUebernehmenDieMaschinen}
		Autor: Klaus Mainzer\\
		Jahr: 2016\\
		ISBN: 978-3-662-48452-4\\
		ISBN: 978-3-662-48453-1 (eBook)

	\subsubsection{Buch: Prinzipien der Mathematik}
	\label{book:prinzipa_of_mathmatics}
		Letzter Zugriff: 24.02.2020 12:14\\
		Autor: Bertrand Russel, Alfred North Whitehead\\
		URL: https://plato.stanford.edu/entries/principia-mathematica/

\subsection{Informationen über das menschliche Gehirn}
\label{subsec:Informationen_ueber_das_menschliche_Gehirn}
	\subsubsection{dasgehirn.info - Zellen: spezialisierte Arbeiter des Gehirns}
	\label{subsubsec:dasgehirn:Zellen-Arbeiter_Des_Gehrins}
	Letzter Zugriff: 22.02.2020 18:03\\
	URL: https://www.dasgehirn.info/grundlagen/kommunikation-der-zellen/zellen-spezialisierte-arbeiter-des-gehirns

\subsection{Viola Jones Algorithmus}
\label{algoryhtm:viola_jones}
	Letzter Zugriff: 02.03.2020 12:52\\
	URL: https://pdfs.semanticscholar.org/40b1/0e330a5511a6a45f42c8b86da222504c717f.pdf

\subsection{Facial Recognition}
\label{source:searchenterpriseai:facial_recognition}
	Letzter Zugriff: 02.03.2020 13:47\\
	URL: https://searchenterpriseai.techtarget.com/definition/facial-recognition

	\subsubsection{Elastic Bunch Graph Matching}
	\label{scholarpedia:ebgm}
	Letzter Zugriff: 02.03.2020 19:42\\
	URL: http://www.scholarpedia.org/article/Elastic\_Bunch\_Graph_Matching


\subsection{\enquote{Machine Learning} und \enquote{Deep Learning}}
\label{subsec:machine_learning_and_deep_learning}
	Letzter Zugriff: 01.03.2020 22:22\\
	URL: https://www.zendesk.com/blog/machine-learning-and-deep-learning/

\subsection{Netflix Research}
\label{Netlfix_Research}
	\subsubsection{Machine Learning - Learn hoe to entertain the World}
	\label{Netflix_Research:machine_learning_learn_how_to_entertain_the_world}
	Letzter Zugriff: 02.03.2020\\
	URL: https://research.netflix.com/research-area/machine-learning

\subsection{Bilder}
\label{subsec:Bilder_Anhang}
	\subsubsection{Beispiel eines Neuronalen Netzes}
	\label{subsubsec:Beispiel_eines_Neuronalen_Netzes}
	Download am: 22.02.2020 15:20\\
	URL: https://de.cleanpng.com/png-jxw2np/


% Bibliografie im .bib Format einfügen
\bibliography{bachelorarbeit}

\end{document}
